\subsection{Genuine native affine point-space gluing}
\label{subsec:native-scheme-atlas}

The gluing here is a topological point-space comparison for actual Mathlib
schemes. The pieces and their section rings are genuine scheme affine opens.
No analytic structure sheaf, coherent-sheaf gluing or Hausdorff assertion for
arbitrary schemes is supplied by this construction.

\begin{theorem}[Actual scheme-open cover gluing]
\label{scheme:native-open-cover-gluing}
\lean{EnvelopeProperty.Berkovich.SchemeValuedPoint.supportOpens,
EnvelopeProperty.Berkovich.SchemeValuedPoint.openCoverGlueData,
EnvelopeProperty.Berkovich.SchemeValuedPoint.OpenCoverGlued,
EnvelopeProperty.Berkovich.SchemeValuedPoint.iUnion_supportOpens,
EnvelopeProperty.Berkovich.SchemeValuedPoint.openCoverGlueHomeomorph}
\leanok
\uses{scheme:valued-open-embedding,scheme:valued-point-topology}
For any family of scheme opens $U_i$ covering $X$, their actual native
support-open inverse images cover the native valued-point space of $X$.
The genuine TopCat gluing of these open subspaces is homeomorphic to that
native space. This statement imposes no finiteness or affineness condition
on the open cover.
\end{theorem}
\begin{proof}
Point projection is continuous, so each scheme open gives an actual native
open. Every underlying scheme support belongs to a member of the given
cover, hence every native valued point belongs to the corresponding native
open. Apply the actual topological open-cover gluing construction to these
opens. Its overlap equivalence identifies the same native point in two
pieces, and its proved open-cover homeomorphism recovers the native space.
The scheme-open restriction homeomorphism identifies each piece with the
actual native valued-point space of that open subscheme.
\end{proof}

\begin{theorem}[Genuine affine section-ring chart parametrization]
\label{scheme:native-affine-atlas}
\lean{EnvelopeProperty.Berkovich.SchemeValuedPoint.affineOpenSpectrumHomeomorph,
EnvelopeProperty.Berkovich.SchemeValuedPoint.affineOpenSpectrumHomeomorph_symm_coe,
EnvelopeProperty.Berkovich.SchemeValuedPoint.affineChartInclusion,
EnvelopeProperty.Berkovich.SchemeValuedPoint.continuous_affineChartInclusion,
EnvelopeProperty.Berkovich.SchemeValuedPoint.openCoverGlueHomeomorph_affineChartInclusion,
EnvelopeProperty.Berkovich.SchemeValuedPoint.isOpenEmbedding_affineChartInclusion,
EnvelopeProperty.Berkovich.SchemeValuedPoint.affineChartInclusion_eq_iff,
EnvelopeProperty.Berkovich.SchemeValuedPoint.affineChartInclusion_jointly_surjective}
\leanok
\uses{scheme:native-open-cover-gluing,scheme:affine-valued-point-homeomorphism,
scheme:valued-open-embedding}
If each $U_i$ is an actual affine open, its support-open native subspace
is canonically homeomorphic to the genuine seminorm spectrum of
$\Gamma(X,U_i)$. These affine spectra give open embeddings into the glued
space and jointly cover it. Their two representatives are identified
exactly when their resulting whole native valued points of $X$ agree.
\end{theorem}
\begin{proof}
Use the actual scheme-open restriction homeomorphism, then the native
homeomorphism induced by the affine scheme isomorphism
$U_i\simeq\operatorname{Spec}\Gamma(X,U_i)$, then the proved native affine
seminorm homeomorphism. This gives the displayed chart homeomorphism. Its
inverse, followed by the actual gluing-piece inclusion, parametrizes each
piece by the spectrum of its genuine section ring. Both inverse maps in
this composition are actual point maps through residue-field isomorphisms;
the ambient representative is the genuine open-inclusion image of that
native point of $U_i$.

Homeomorphisms followed by actual gluing-piece open embeddings are open
embeddings, hence continuous. Joint surjectivity of the genuine gluing
inclusions and surjectivity of the chart homeomorphisms prove coverage.
The actual open-cover gluing homeomorphism sends a chart representative
to precisely its ambient native point. Its injectivity proves both directions
of the representative equality criterion. This does not identify an analytic
structure sheaf or assert a scheme/sheaf gluing theorem.
\end{proof}

\begin{theorem}[Canonical atlas of every actual scheme]
\label{scheme:canonical-native-affine-atlas}
\lean{EnvelopeProperty.Berkovich.SchemeValuedPoint.AffineAtlasIndex,
EnvelopeProperty.Berkovich.SchemeValuedPoint.canonicalAffineOpens,
EnvelopeProperty.Berkovich.SchemeValuedPoint.isAffineOpen_canonicalAffineOpens,
EnvelopeProperty.Berkovich.SchemeValuedPoint.iUnion_canonicalAffineOpens,
EnvelopeProperty.Berkovich.SchemeValuedPoint.canonicalAffineGlueHomeomorph}
\leanok
\uses{scheme:native-affine-atlas,scheme:native-open-cover-gluing}
Index by all actual affine opens of $X$. This is a genuine covering family,
so its native affine section-ring point-space gluing recovers the valued-point
space of every scheme $X$.
\end{theorem}
\begin{proof}
Every actual scheme point has an affine open neighborhood by the scheme's
proved affine basis. Such an open, with its actual affineness proof, is one of
the indexing objects; hence the union of the indexed opens is all of $X$.
Apply the arbitrary-cover native gluing homeomorphism and the affine chart
parametrizations to this concrete family. This provides an actual canonical
atlas without assuming that a requested cover exists or is finite.
\end{proof}
